\typeout{IJCAI--ECAI 26 Instructions for Authors}
\documentclass{article}
\pdfpagewidth=8.5in
\pdfpageheight=11in
\usepackage{ijcai26}
\newif\ifanonymized
\anonymizedtrue
\newif\ifreview
\reviewtrue
\usepackage{times}
\usepackage{soul}
\usepackage{url}
\usepackage{xurl}
\usepackage[hidelinks]{hyperref}
\usepackage[utf8]{inputenc}
\usepackage[small]{caption}
\usepackage{graphicx}
\usepackage{amsmath}
\usepackage{amsthm}
\usepackage{amssymb}
\usepackage{booktabs}
\usepackage{multirow}
\usepackage{enumitem}
\usepackage{algorithm}
\usepackage{algorithmic}
\usepackage[switch]{lineno}
\ifreview
\linenumbers
\fi
\urlstyle{same}
\newtheorem{example}{Example}
\newtheorem{theorem}{Theorem}
\pdfinfo{
/TemplateVersion (IJCAI.2026.0)
}

\title{OrthKD: Extracting Generalized Clinical Knowledge from Heterogeneous Teachers for Lightweight Deployment}

\ifanonymized
\author{Anonymous Author(s)}
\else
\author{
Yi Xu$^{1,*}$
\and
Cheng Chen$^{1,*}$
\and
Mufan Cao$^1$
\affiliations
$^1$Tongji University, Shanghai, China
\emails
\{2351441, 2351445, 2352738\}@tongji.edu.cn
}
\fi


\begin{document}

\maketitle

\begin{abstract}
Primary healthcare requires edge-efficient models with robust generalization across diverse populations. Multi-teacher knowledge distillation (KD) can compress ensembles into lightweight students, but faces a critical challenge: expertise asymmetry. In diabetic retinopathy (DR) screening, a Transformer teacher (Swin-Base: 0.83 QWK) significantly lags behind a strong CNN teacher (EfficientNet-B3: 0.88 QWK). Standard fusion fails here as erroneous logits from the weak teacher bias the student. Therefore, we propose OrthKD (Orthogonal Feature Knowledge Distillation): extracting valuable global context from the weak teacher only at the feature layer while imposing orthogonality constraints to ensure complementarity, thereby circumventing unreliable logit supervision. Experiments on 132,049 retinal images confirm our lightweight MobileNetV3 student (5.4M parameters) achieves strong in-domain performance (0.885 QWK on EyePACS) and significantly improves zero-shot generalization on Messidor-2 (0.728 QWK vs. 0.507 baseline), enabling robust DR screening on resource-constrained devices.
\end{abstract}

\section{Introduction}

\subsection{Technical Challenges in Medical AI Deployment}
Deep learning models now rival ophthalmologists in diabetic retinopathy (DR) screening~\cite{gulshan2016development,zhou2021review}, yet clinical impact remains limited. Haq et al.~\cite{haq2024computationally} analyzed 84 DR AI papers, noting that neglected computational complexity is a major shortcoming. While large models (EfficientNet-B7, ResNet-152, ViT-L) may achieve relatively high performance, they require GPU infrastructure and cloud connectivity that are difficult for primary healthcare settings to access. Meanwhile, 537 million high-risk individuals globally still lack screening accessibility.~\cite{teo2021global}.

Some lightweight architectures (MobileNetV3, EfficientNet-Lite, ShuffleNetV2)~\cite{howard2017mobilenets,zhang2018shufflenet} support CPU inference but sacrifice sensitivity at adjacent classification boundaries—precisely the regions where annotator agreement is lowest (Cohen's kappa=0.61–0.68).~\cite{krause2018grader}. This performance degradation is particularly detrimental to the DR screening we focus on because: (1) Sequential structure is critical—DR severity progresses monotonically through grades 0–4, with referable-DR threshold at grade $\geq$2 defining treatment urgency; (2) Ambiguity at boundaries is objectively present—the transition between grades 1–2–3 accounts for 68\% of expert disagreement; (3) Asymmetric costs of referral decisions—missed diagnoses at the referral threshold delay treatment for vision-threatening diseases. Therefore, our primary goal is to ensure peak diagnostic performance while enabling deployment-friendly lightweight models with stronger cross-domain generalization capabilities to meet real-world clinical implementation needs.

\subsection{Limitations of Standard Knowledge Distillation}
Knowledge distillation (KD)~\cite{hinton2015distilling} compresses teacher knowledge into lightweight student models by matching probability distributions. However, when fusing heterogeneous architectures (CNN + Transformer), significant capability gaps may emerge between teachers, ultimately leading to suboptimal model performance. This paper proposes a novel approach to address this expertise asymmetry through selective knowledge transfer.

This limitation is particularly consequential in healthcare deployment. First, real-world screening sites rarely match the training distribution due to differences in cameras, acquisition protocols, and patient populations. Second, clinical decision-making is often threshold-driven (e.g., referable DR defined by grade $\geq 2$), where label ambiguity and decision costs are higher. Therefore, it is not sufficient for distillation to improve in-domain accuracy: the compressed model must also remain reliable and robust under distribution shift. This motivates our emphasis on zero-shot cross-dataset evaluation and the use of safety- and reliability-oriented metrics (AUC and calibration) alongside ordinal grading accuracy.

\subsection{The Expertise Asymmetry Challenge}
We made three key discoveries during our experiments.

First, expertise asymmetry: The Transformer teacher (Swin-Base, 0.83 QWK) significantly lags behind both the CNN teacher (EfficientNet-B3, 0.88 QWK) and the student baseline (0.85 QWK). This expertise asymmetry is common when teachers exhibit substantial differences in inductive bias or capacity.

Second is the complementary knowledge possessed by weak teachers. Despite its overall low accuracy, the Transformer correctly classified 2.6\% of gold samples—cases misclassified by both the CNN and student models. These samples typically rely on global contextual information (e.g., subtle retinal thinning) rather than local, fine-grained lesions.

Third, logit fusion amplifies harm. Naive logit distillation on teachers with expertise asymmetry achieves only 73\% accuracy on “disputed” samples—weak teachers introduce a 1:6 harm-to-benefit ratio, confining the student model within the teacher's systemic bias.

We contend that effective distillation should capture weak teachers' unique expertise while circumventing their systemic errors. Our efforts yielded promising results: OrthKD achieved a +22.1 percentage point improvement in out-of-distribution (OOD) generalization on external datasets compared to the baseline.

\subsection{OrthKD: Our Solution}
We propose OrthKD (Orthogonal Feature Knowledge Distillation), which addresses expertise asymmetry through two key innovations. First, we employ \textit{feature-only distillation from weak teachers}: Transformers provide only features, not logit supervision, thereby extracting valuable global context while avoiding classification errors that could contaminate the learning process. Second, we impose \textit{orthogonal constraints for generalization}: we penalize redundancy between the feature projections from CNN and ViT, serving as a diversity regularization that compels students to learn more transferable complementary patterns.

\subsection{Contributions}
This paper makes four main contributions:
\begin{itemize}[leftmargin=*,itemsep=2pt,topsep=2pt]
    \item \textit{Deployment efficiency:} Enables the lightweight MobileNetV3 student (5.4 million parameters) to rival ``heavyweight'' heterogeneous teacher combinations (totaling over 100 million parameters), making high-quality DR screening feasible on edge devices without GPUs.
    \item \textit{Solutions for expertise asymmetry:} Introducing OrthKD, which systematically extracts usable value from weak teachers by decoupling feature learning from unreliable logit supervision.
    \item \textit{Improved generalization capability:} OrthKD achieved a +22.1 percentage point absolute improvement over QWK on unseen distributions (Messidor-2: 0.728 vs 0.507 baseline), bridging the gap between academic benchmarks and real-world clinical utility.
    \item \textit{Counterintuitive insight:} Demonstrates that even a subpar Transformer (0.83 QWK) can contribute critical global context when properly distilled, challenging the conventional wisdom that ``stronger teachers are always better.''
\end{itemize}

\section{Related Work}

Knowledge distillation (KD)~\cite{hinton2015distilling} transfers knowledge by minimizing KL divergence between teacher and student logits. Multi-teacher KD~\cite{you2017learning,mirzadeh2020improved} typically assumes that all teachers provide equivalent high-quality guidance. Recent heterogeneous distillation works~\cite{passalis2020heterogeneous} address the architecture mismatch (CNN to Transformer), but still assume comparable teacher performance. In contrast, our work identifies and addresses the Expertise Asymmetry scenario where a weak teacher (e.g., underperforming ViT) is leveraged for its unique complementary features despite its unreliable decisions.

Feature-level distillation~\cite{romero2014fitnets} and attention transfer~\cite{zagoruyko2016paying} enable students to learn richer intermediate representations. However, naive multi-teacher feature alignment may lead to redundant knowledge acquisition, impairing generalization performance. Diversity regularization (often achieved through variance pooling or adversarial training) aims to learn decoupled patterns. OrthKD addresses this by imposing strict orthogonality constraints on teacher-specific projection heads, ensuring students extract distinct and non-overlapping semantic information from each teacher.

Deploying AI for DR screening in resource-constrained primary care settings requires both efficiency and cross-dataset robustness~\cite{cleland2023artificial,zhou2021review}. While standard compression techniques (quantization, pruning) yield lightweight models~\cite{howard2017mobilenets,han2015deep}, they often overfit to the training distribution and fail on unseen clinical populations. Our work demonstrates that combining orthogonally driven dual-teacher KD with edge optimization architecture (MobileNetV3) enhances out-of-distribution generalization (+22.1 pp QWK on Messidor-2) while providing necessary diversity.

Traditional methods often enforce strict spatial alignment between student and teacher attention maps~\cite{zagoruyko2016paying}. However, in our asymmetric expertise setup, weak teacher attention captures global context but lacks the local fineness of CNNs; strict alignment forces students to learn these shortcomings, thereby capping performance below the ceiling of strong teachers. In contrast, OrthKD's orthogonal constraints foster complementarity: students can retain CNN precision while selectively absorbing Transformer-based global context.

% =========================================================================
% SECTION 3: METHODOLOGY
% =========================================================================
\section{Methodology}

OrthKD addresses expertise asymmetry through two core innovations: (1) Feature Distillation Only: Extracting global context from weak Transformer teachers while bypassing unreliable logits; (2) Orthogonal Constraints: Ensuring knowledge transfer is complementary rather than redundant. The framework is illustrated in the Figure~\ref{fig:framework}.

\paragraph{Feature representation.}
Let $z$ denote pre-softmax logits and let $f$ denote the intermediate representation used for distillation. In all networks, we take $f$ from the final backbone stage before the classification head. If the native representation is spatial, we apply global average pooling to obtain a single vector. This yields compact features $f^S, f^{CNN}, f^{ViT} \in \mathbb{R}^{d}$, which are compatible with lightweight deployment and stable feature alignment.

\begin{figure*}[t]
    \centering
    \includegraphics[width=0.95\textwidth]{orthkd_framework.pdf}
    \caption{\textbf{OrthKD Framework.}
    Given a fundus image, three networks extract features in parallel: a strong CNN teacher (EfficientNet-B3, 0.88 QWK), a lightweight student (MobileNetV3), and a weak ViT teacher (Swin-Base, 0.83 QWK).
    The student learns through two projection heads: $g_{cnn}$ aligns with CNN features and receives \textit{full supervision} (logits + features), while $g_{vit}$ aligns with ViT features but receives \textit{feature-only supervision} (ViT logits are blocked).
    The orthogonality constraint $\mathcal{L}_{orth}$ enforces complementary projected features, discouraging redundancy.}
    \label{fig:framework}
\end{figure*}

\subsection{Problem Formulation}

We address five-grade DR severity classification $y \in \{0,1,2,3,4\}$ with ordinal structure following the ICDR scale~\cite{wilkinson2003proposed}. We select EfficientNet-B3 ($\mathcal{T}_{CNN}$, 0.88 QWK) as the strong CNN teacher and Swin-Base ($\mathcal{T}_{ViT}$, 0.83 QWK) as the Transformer teacher---the most competitive among tested ViT variants. This asymmetry in expertise (CNN $>$ ViT) is not merely an experimental anomaly, but a common characteristic in medical imaging: unlike natural vision where ViTs thrive on massive datasets (e.g., JFT-300M), medical domains are data-scarce relative to model capacity. Under this regime, CNNs' inductive biases (e.g., translation invariance) provide a persistent advantage.

The student is MobileNetV3-Large (5.4M params), achieving 0.853 QWK when trained independently. Standard dual-teacher distillation that combines logits from both teachers leads to performance degradation---the weaker Transformer introduces erroneous guidance. Our analysis shows only 73\% accuracy on controversial samples, with a harm-to-benefit ratio of 1:6 (Section~\ref{sec:analysis}).
\label{sec:analysis}

\subsection{Orthogonal Feature Distillation}
Although the weak Transformer teacher has lower classification accuracy, it can capture unique global contextual features---it correctly classifies 2.6\% of the samples where both the CNN and the student model fail. We extract this knowledge only at the feature layer, while bypassing the unreliable logits. Specifically, we decompose knowledge transfer as follows:
\begin{itemize}[leftmargin=*,itemsep=1pt,topsep=2pt]
    \item \textit{CNN Teacher (Strong):} Full transfer---logits + features.
    \item \textit{ViT Teacher (Weak):} Feature-only transfer---no logit supervision.
\end{itemize}


\begin{figure}[t]
    \centering
    \includegraphics[width=\columnwidth]{vis_gradcam.pdf}
    \caption{\textbf{Grad-CAM visualization.}
    CNN focuses on local textures (e.g., optic disc), while ViT exhibits diffuse attention capturing global context.
    The baseline student overfits to salient landmarks, whereas the OrthKD student localizes lesions and better captures vascular structures.}
    \label{fig:gradcam}
\end{figure}

\subsection{Qualitative Analysis}
To intuitively understand how orthogonality improves feature representation, we visualize Grad-CAM in Figure~\ref{fig:gradcam}.

As shown in the second row (normal sample), the CNN teacher and baseline student primarily focus on the optic disc, a highly prominent local landmark. This indicates that without proper guidance, the lightweight student tends to "learn shortcuts"—classifying based on simple landmarks, especially bright areas, rather than understanding the retinal structure. In contrast, the ViT teacher has a more global receptive field and focuses on the entire retinal background to confirm the absence of lesions.

More importantly, the OrthKD student model (column 5) exhibits a unique "hybrid" attention pattern: while retaining the sharpness typical of CNN structures (derived from the convolutional architecture), its activation regions are significantly larger compared to the baseline. Specifically, it successfully tracks the arc-shaped trajectory of blood vessels extending from the optic disc, a behavior absent in the baseline but consistent with the global perception of ViT.

This validates our hypothesis: by imposing orthogonality between student projection features, OrthKD avoids model collapse to solutions that only mimic the easiest features (optic disc); instead, it forces the student to explore orthogonal, complementary subspaces, thereby learning richer semantic information (such as vascular topology), achieving a combination of CNN precision and ViT breadth.


Let $f^S, f^{CNN}, f^{ViT} \in \mathbb{R}^d$ be student and teacher feature representations. The Orthogonal Feature KD loss combines three components: (1) logit KD from the CNN teacher,
\begin{equation}
\mathcal{L}_{\mathrm{KD}}^{CNN} = \tau^2 \cdot \mathrm{KL}\left(\sigma(z^S/\tau), \sigma(z^{CNN}/\tau)\right)
\end{equation}
where $\sigma(\cdot)$ is the softmax and $\tau$ is the distillation temperature; (2) feature alignment with both teachers,
\begin{equation}
\begin{aligned}
\mathcal{L}_{\text{feat}}^{CNN} &= \|g_{cnn}(f^S) - f^{CNN}\|_2^2, \\
\mathcal{L}_{\text{feat}}^{ViT} &= 1 - \cos(g_{vit}(f^S), f^{ViT})
\end{aligned}
\end{equation}
where $g_{cnn}(\cdot)$ and $g_{vit}(\cdot)$ are learnable projection layers adapting student features to each teacher's latent space; and (3) the orthogonality constraint,
\begin{equation}
\mathcal{L}_{\text{orth}} = \cos^2(g_{cnn}(f^S), g_{vit}(f^S))
\label{eq:orth}
\end{equation}
 This penalizes the correlation between the student's two projection heads. We adopt the squared cosine similarity as it provides smoother gradients near zero compared to absolute-value formulations, facilitating more stable optimization of distinct feature subspaces for satisfying both teachers simultaneously.

\subsection{Training Objective}
The complete OrthKD loss combines classification, knowledge distillation, and orthogonality:
\begin{equation}
\mathcal{L}_{\mathrm{total}} = \alpha_1 \mathcal{L}_{\mathrm{CE}} + \alpha_2 \mathcal{L}_{\mathrm{KD}}^{CNN} + \alpha_3 \mathcal{L}_{\text{feat}}^{CNN} + \alpha_4 \mathcal{L}_{\text{feat}}^{ViT} + \alpha_5 \mathcal{L}_{\text{orth}}
\end{equation}
where $\mathcal{L}_{\mathrm{CE}}$ is cross-entropy with ground truth, $\mathcal{L}_{\mathrm{KD}}^{CNN}$ is logit distillation from the CNN teacher, $\mathcal{L}_{\text{feat}}$ are feature alignment losses, and $\mathcal{L}_{\text{orth}}$ enforces orthogonality. Hyperparameters: $(\alpha_1, \alpha_2, \alpha_3, \alpha_4, \alpha_5) = (0.2, 0.4, 0.2, 0.1, 0.1)$.


Images are resized to 512$\times$512 with augmentation including random flips, rotations, CLAHE, color jitter, and Cutout. We use AdamW optimizer (lr=3e-4, weight decay=1e-4) with cosine annealing for 50 epochs (batch size 32) on NVIDIA GPUs. Loss weights $\alpha_1=0.2, \alpha_2=0.4, \alpha_3=0.2, \alpha_4=0.1, \alpha_5=0.1$ are fixed across all experiments to ensure rigorous evaluation of generalization capability.

\subsection{Why Orthogonality Improves OOD Robustness}
A significant observation in our experiments is that orthogonal constraints have almost negligible impact on in-domain performance but lead to significant improvements in out-of-domain performance. Specifically, removing orthogonality only causes a slight degradation on the training distribution (Table~\ref{tab:ablation}), but leads to a noticeable decline on external clinical datasets such as Messidor-2. We attribute this to three complementary mechanisms.

First, orthogonality prevents \textit{redundancy collapse}. Without explicit decorrelation, the two projection heads $g_{cnn}$ and $g_{vit}$ tend to converge to similar representations---encoding the ``easiest'' discriminative features (e.g., the optic disc). Minimizing $\cos^2(g_{cnn}(f^S), g_{vit}(f^S))$ forces them into different subspaces, allowing the student to simultaneously capture local lesion evidence from the CNN and global structural context from the ViT.

Second, orthogonality \textit{mitigates shortcut learning}. Lightweight students tend to exploit spurious correlations associated with salient landmarks. The weaker Transformer, despite lower accuracy, focuses on background structures less dependent on domain-specific features. Orthogonality constraints force the student to retain this ``background-aware channel,'' regularizing against shortcut reliance.

Third, orthogonality provides \textit{diversity as a buffer against distribution shift}. When training and external distributions differ (e.g., due to imaging equipment or population characteristics), redundant features degrade simultaneously. Orthogonal subspaces provide decorrelated evidence streams; when one stream degrades, the other may remain stable, yielding more robust predictions.

In summary, orthogonality $\to$ complementary feature subspaces $\to$ reduced overfitting to domain-specific shortcuts $\to$ stronger OOD generalization. This causal chain explains why orthogonality is ``silent'' within a single domain but crucial in real-world clinical deployments.

\subsection{Deployment Considerations}
We select MobileNetV3-Large (5.4M parameters) as the student architecture for its deployment-friendly design:
\begin{itemize}[leftmargin=*,itemsep=2pt]
\item \textit{Parameter efficiency:} 5.4M params (vs. 12M B3, 88M Swin).
\item \textit{Mobile-optimized:} Inverted residuals and hard-swish activations.
\item \textit{Edge deployment:} Designed for CPU inference without GPU requirements.
\end{itemize}
This makes DR screening possible in resource-constrained primary healthcare settings that lack GPU infrastructure.

Beyond efficiency, clinical deployment also requires models to perform stably near the referral threshold (grade $\geq$ 2) and be robust to variations in image quality. Our strong performance on Messidor-2 (0.728 QWK)—a dataset that likely exhibits greater domain shift and variability than the high-quality IDRiD dataset—validates the potential of OrthKD in primary care settings. The OrthKD student model inherits the clinical accuracy of CNNs while incorporating the complementary perspective of Transformers, resulting in more stable predictions at the decision boundary where clinical risk is highest. In a typical primary care triage process, this model can assist non-specialist screening personnel in flagging referral cases for ophthalmologist review, reducing workload while maintaining a safety margin.
\section{Experiments}

\subsection{Datasets and Experimental Protocol}
We combine two public DR screening datasets with data augmentation:
\begin{itemize}[leftmargin=*,itemsep=1pt,topsep=1pt]
    \item \textit{EyePACS} (Kaggle Diabetic Retinopathy Detection Challenge, 2015): 126,584 fundus images from California screening sites, labeled with 5-grade ICDR scales~\cite{wilkinson2003proposed}.
    \item \textit{APTOS} (Asia Pacific Tele-Ophthalmology Society Challenge, 2019): 5,465 images from Indian population, using the same 5-grade labeling~\cite{aptos2019}.
\end{itemize}
Total training pool: 132,049 base images, expanded via comprehensive data augmentation. We use a stratified split: 80\% train, 10\% validation, 10\% test. Teachers and student trained on the combined dataset.

To assess cross-dataset generalization under distribution shift:
\begin{itemize}[leftmargin=*,itemsep=1pt,topsep=1pt]
    \item \textit{Messidor-2} (European population): N=1,728, macula-centered images with expert adjudication. Not seen during training.
    \item \textit{IDRiD} (Indian Diabetic Retinopathy Image Dataset)~\cite{porwal2018indian}: N=516, high-resolution fundus images from an Indian clinical setting with expert-verified 5-level grading. This dataset represents real-world clinical validation with different population demographics and imaging equipment.
\end{itemize}

The black borders were removed using circular cropping, and CLAHE (clip limit=2.0) was used for illumination normalization, followed by scaling to 512$\times$512. No specific parameter tuning was performed on the external dataset (zero-shot evaluation).

The main reporting metric in this paper is the quadratic weighted Kappa (QWK), which measures the consistency of the five-level grading system. To evaluate the safety and reliability of clinical screening, we additionally report the ROC-AUC and Expected Calibration Error (ECE, 15 bins) for referable DR (grade $\geq 2$) on the in-domain EyePACS and APTOS datasets. For external validation sets (Messidor-2 and IDRiD), given the characteristics of zero-shot testing, we focus on QWK. This is because QWK most directly reflects the model's inherent generalization ability in ordinal ranking of disease severity without requiring threshold calibration for the target domain, thus avoiding evaluation bias caused by differences in threshold selection.

\subsection{Implementation Details}
We train teacher models and the student baseline on the combined training pool, and then train student variants under different distillation configurations while keeping the student architecture fixed (MobileNetV3-Large). For fair comparison, we keep preprocessing and augmentation identical across all methods, and we avoid dataset-specific tuning on external test sets to reflect realistic deployment conditions. Unless otherwise specified, we keep key training hyperparameters and loss weights consistent across student variants to isolate the effect of distillation design.

We compare our method against five configurations: (1) \textit{EfficientNet-B3}, the CNN teacher (12M params, 0.880 QWK); (2) \textit{Swin-Base}, the ViT teacher (88M params, 0.830 QWK) introducing structural diversity; (3) \textit{MobileNetV3}, the student baseline (5.4M params, 0.853 QWK); (4) \textit{B3-Only KD}, standard single-teacher distillation using only CNN cues; (5) \textit{CW Logit Fusion}, a multi-teacher baseline using confidence-weighted logits; and (6) \textit{OrthKD (Ours)}, the proposed orthogonal feature distillation from heterogeneous teachers.

% =========================================================================
% SECTION 5: RESULTS
% =========================================================================
\section{Results}

\subsection{Main Results: Clinical Performance and Safety}

In actual clinical diagnosis, the three most frequently used metrics are grading accuracy (QWK), screening safety (AUC for referral vs. non-referral), and reliability (Expected Calibration Error, ECE).

Table~\ref{tab:main_results} provides a comprehensive evaluation of the model's performance based on the three metrics mentioned above.

OrthKD achieved the highest classification accuracy on two domain datasets (EyePACS: 0.885 QWK; APTOS: 0.952). More importantly, it performed comparably to or even better than strong teacher models in terms of screening safety, achieving the highest AUC (0.966 / 0.995). This demonstrates that model compression does not weaken its ability to identify referral cases, which is a core requirement for clinical safety.

Additionally, OrthKD demonstrates superior calibration (ECE = 0.023 on EyePACS), significantly outperforming the overconfident CNN teacher (0.045) and the baseline student (0.037). Comparing ``w/o Orthogonality'' and ``OrthKD,'' although the QWK scores are similar, we observe a significant decrease in ECE ($0.035 \to 0.023$), indicating that the orthogonality constraint effectively regularizes the student's confidence, preventing it from collapsing into a redundant or overconfident feature space.

\textit{Note on Swin Transformer:} Swin-Base achieves an ECE close to 0 (0.005) on EyePACS, but its accuracy is lower (0.830 QWK). This suggests that a low ECE might reflect underfitting (overall low confidence) rather than good calibration at high accuracy. In contrast, OrthKD achieves a low ECE while maintaining high accuracy, demonstrating better reliability.

\begin{table*}[t]
\centering
\caption{\textbf{Comprehensive In-Domain Performance.} We evaluate models on \textbf{QWK} (Accuracy), \textbf{AUC} (Safety), and \textbf{ECE} (Reliability). OrthKD achieves the best balance: matching the teacher's safety (AUC) while significantly improving calibration (ECE) and grading precision (QWK).}
\label{tab:main_results}
\small
\setlength{\tabcolsep}{5pt}
\begin{tabular}{@{}lcccccc@{}}
\toprule
\multirow{2}{*}{\textbf{Method}} & \multicolumn{3}{c}{\textbf{EyePACS (In-Domain)}} & \multicolumn{3}{c}{\textbf{APTOS (In-Domain)}} \\
\cmidrule(lr){2-4} \cmidrule(l){5-7}
 & \textbf{QWK} $\uparrow$ & \textbf{AUC} $\uparrow$ & \textbf{ECE} $\downarrow$ & \textbf{QWK} $\uparrow$ & \textbf{AUC} $\uparrow$ & \textbf{ECE} $\downarrow$ \\
\midrule
Swin-Base (ViT Teacher) & .830 & .935 & \textbf{.005}$^\dagger$ & .918 & .980 & .038 \\
EffNet-B3 (CNN Teacher) & .876 & .965 & .045 & .941 & .993 & .032 \\
MobileNetV3 (Baseline) & .853 & .947 & .037 & .934 & .987 & .030 \\
\midrule
B3-Only KD & .878 & .965 & .027 & .942 & .994 & .040 \\
w/o Orthogonality & .883 & .965 & .035 & .950 & .995 & .031 \\
\textbf{Ours (OrthKD)} & \textbf{.885} & \textbf{.966} & .023 & \textbf{.952} & \textbf{.995} & \textbf{.018} \\
\bottomrule
\multicolumn{7}{l}{\scriptsize $^\dagger$ Swin's low ECE on EyePACS may reflect underfitting (overly low-confidence predictions) rather than well-calibrated high-accuracy predictions.}
\end{tabular}
\end{table*}

\subsection{Ablation Study}
\label{sec:ablation}
Table~\ref{tab:ablation} verifies the contribution of each component. Removing the ViT features (B3-Only) results in a 0.9\% performance drop, indicating that the weak teacher indeed provides valuable complementary knowledge. Orthogonal constraints have only a marginal impact within the domain (approximately 0), while confidence-weighted logit fusion performs significantly worse (-1.6\%), further confirming that using weak teacher logits should be avoided.

\begin{table}[ht]
\centering
\caption{\textbf{Ablation Study.} Each row removes one component from the full OrthKD.}
\label{tab:ablation}
\small
\begin{tabular}{@{}lccc@{}}
\toprule
\textbf{Variant} & \textbf{APTOS} & \textbf{EyePACS} & \textbf{$\Delta$} \\
\midrule
Full OrthKD & .952 & .885 & — \\
	w/o ViT Features (B3-Only) & .942 & .878 & -0.9\% \\
	w/o Orthogonality & .950 & .883 & $\approx$0 \\
	CW Logit Fusion & .937 & .869 & -1.6\% \\
\bottomrule
\end{tabular}
\end{table}

\subsection{Cross-Dataset Generalization}

Our core hypothesis is that orthogonality primarily acts as a diversity regularizer under distribution shift, rather than improving in-domain performance. Table~\ref{tab:crossdataset} supports this: although orthogonality provides only a small gain on the training set (see Table~\ref{tab:ablation}), it brings a +7.4 percentage point improvement on Messidor-2 (0.728 vs 0.654), significantly narrowing the gap with the strong teacher model.

In external validation, we emphasized zero-shot generalization for realism, and therefore did not employ dataset-specific threshold tuning or post-processing calibration. Consequently, we used the Quadratic Weighted Kappa (QWK) on Messidor-2 and IDRiD as the primary metric, as it directly measures ordinal classification agreement and is insensitive to operating point selection during deployment. Within the domain, we additionally report referral AUC and ECE to characterize screening safety and reliability under matching distributions.

\begin{table}[ht]
\centering
\caption{\textbf{External Validation (Zero-Shot).} Comparison of all methods on unseen clinical datasets.}
\label{tab:crossdataset}
\small
\begin{tabular}{@{}lcc@{}}
\toprule
\textbf{Method} & \textbf{Messidor-2} & \textbf{IDRiD} \\
\midrule
B3 (CNN Teacher) & .785 & .749 \\
Swin (ViT Teacher) & .453 & .636 \\
MobileNetV3 (Baseline) & .507 & .516 \\
\midrule
B3-Only KD & .587 & .701 \\
No Orthogonality & .654 & .724 \\
CW Logit Fusion & .659 & .728 \\
\textbf{Ours (OrthKD)} & \textbf{.728} & \textbf{.743} \\
\bottomrule
\end{tabular}
\end{table}

\begin{figure}[t]
    \centering
    \includegraphics[width=\columnwidth]{evidence_case.pdf}
    \caption{\textbf{Ambiguity resolution case study.}
    On IDRiD\_303 (Grade 2), the baseline student (B3-Only) overfits to local textures and predicts Grade 0, while the CNN teacher predicts Grade 2.
    Although the weak ViT's logits are unreliable, its features provide global structural context; OrthKD leverages this feature-level complementarity via orthogonality.}
    \label{fig:ambiguity}
\end{figure}

OrthKD outperforms all other student variants across both external datasets. Compared to the single-teacher B3-Only model, it improves by 14.1 percentage points on Messidor-2 and 4.2 percentage points on IDRiD. This suggests that capturing the ``global intelligence'' of weak Transformers is more effective than simply mimicking the most accurate teacher. Notably, feature transfer from Swin (``No Orthogonality'') already outperforms B3-Only KD on Messidor-2 (0.654 vs. 0.587), supporting our claim that weak-teacher features contain useful complementary evidence when decoupled from unreliable logits.

% =========================================================================
% SECTION 6: DISCUSSION
% =========================================================================
\section{Discussion}
\subsection{Key Insight: Orthogonality Enables Generalization}
Our experiments show that orthogonal constraints are a key driving factor for generalization gains, rather than improving in-domain performance. On the training dataset (APTOS/EyePACS), OrthKD showed only minimal improvement compared to the non-orthogonal approach (approximately 0\% in Table~\ref{tab:ablation}). However, this difference widened significantly on external validation: a +7.4 percentage point improvement on Messidor-2 (0.728 vs 0.654).

This demonstrates that orthogonality acts as a diversity regularizer here—forcing students to learn complementary patterns from heterogeneous teachers, rather than learning redundant features that would overfit the training distribution. This is key to improving generalization performance.

\subsection{AI \& Health Perspective: Screening Safety and Deployment}
In an AI \& Health screening workflow, model utility depends not only on ordinal grading accuracy, but also on safe behavior around clinically meaningful thresholds (e.g., referable DR) and robustness to site-to-site variability. This is why we report AUC and calibration (ECE) in addition to QWK, and emphasize zero-shot external validation without target-domain tuning. OrthKD is designed for such settings: it retains the strong teacher's decision-level supervision, while using the weak teacher only to inject complementary global context in a controlled feature-level manner.

\subsection{The Discovery Journey}
After several failures and analyses, we arrived at OrthKD.

\textit{Attempt 1: Confidence-weighted logit fusion.} We first tried weighted fusion of teacher logits based on prediction confidence. The results were poor; high-confidence errors from the weak Transformer led to a harm-to-benefit ratio of 1:6 (73\% accuracy).

\textit{Attempt 2: Feature distillation only.} We hypothesized that Swin's logits were unreliable but its features could capture useful global context. We had the CNN teacher provide logits + features while Swin provided features only. This yielded in-domain improvement (+1\%), but limited generalization gains.

\textit{Attempt 3: Orthogonality.} Analysis revealed that CNN and Swin features were learning redundant patterns---the cosine similarity was very high. Penalizing feature correlation resulted in significant generalization improvement (Messidor-2 +7.4 pp).

\subsection{Why Single-Teacher KD Generalizes Poorly}
An unexpected yet understandable finding is that our strongest teacher B3 only achieved 75\% of its generalization capability under single-teacher KD (0.587 vs 0.785 on Messidor-2). We speculate two mechanisms are at play: (1) \textit{capacity bottleneck}---a 5.4M-parameter student cannot perfectly fit a 12M-parameter teacher, and imperfect distillation amplifies domain-specific patterns; (2) \textit{lack of diversity}---single-teacher KD lacks the regularization effect of learning from multiple perspectives. This explains why OrthKD significantly outperforms B3-Only: the weaker teacher (Swin) contributes crucial diversity in direction and logic to the model's learning.

\subsection{Implications beyond DR}
Although this paper focuses on DR grading, we believe OrthKD principles apply broadly. In any scenario where heterogeneous teachers exhibit significant capability differences---such as chest X-ray analysis where a CNN excels at local nodule detection while a Transformer captures global lung field symmetry---OrthKD provides a systematic method for leveraging diverse and unequal expert models. This enables robust and reliable clinical AI free from unreliable decision-level signals, reducing burden on physicians.

\subsection{Limitations}
This work has two main limitations. First, regarding \textit{scalability}, this paper only tested two teachers; further research should explore extensions to more teachers with pairwise orthogonality constraints. Second, regarding \textit{theoretical understanding}, why orthogonality significantly enhances generalization but minimally impacts in-domain performance remains an open question requiring further investigation.


\section{Conclusion}

Faced with the demand for DR screening on edge devices, OrthKD successfully addressed a long-neglected problem: how to extract maximum value from a heterogeneous ensemble of ``specialized'' teachers. This paper not only proposes an efficient model with only 5.4M parameters, but more importantly, it validates a new heterogeneous distillation paradigm that completely decouples feature learning from unreliable logit supervision.

By imposing orthogonality constraints, we encourage the student to learn feature subspaces complementary to those of the strong CNN teacher, achieving a 22.1 percentage point improvement in generalization performance on the unseen Messidor-2 dataset. Overall, OrthKD provides a practical heterogeneous distillation paradigm that decouples weak-teacher feature transfer from unreliable logit supervision, enabling robust DR screening on edge devices.

\section*{Declarations}

During the preparation of this work, the authors used large language models (Gemini) solely for grammatical correction and polishing of the text. All scientific claims, experimental results, and logical deductions are original contributions of the authors.

\bibliographystyle{named}
\bibliography{ijcai26}

\end{document}
